% Ochrelune PDF 版式（XeLaTeX）。用法：正文第一行 % Ochrelune PDF 版式（XeLaTeX）。用法：正文第一行 % Ochrelune PDF 版式（XeLaTeX）。用法：正文第一行 % Ochrelune PDF 版式（XeLaTeX）。用法：正文第一行 \input{ochrelune-pdf.tex}
% 需要在主文件之前（或主文件第一行之前）定义好：
%   \def\ochretitle{文档标题}   \def\ochrefooter{主题 · 日期}（页脚不写署名）
% 打印版：再加一行 \def\ochreprint{}，纸面和卡片底改为纯白（用户自备米色纸，黑白打印机印底色会出淡墨）。
% 字体放在同目录 fonts/ 下（assets/fetch-fonts.sh 会下载）。
\documentclass[11pt]{article}
\usepackage[a4paper,top=22mm,bottom=24mm,left=22mm,right=22mm,headsep=8mm,footskip=12mm]{geometry}
\usepackage{fontspec}
\usepackage[CJKmath=true]{xeCJK}
\usepackage{xcolor}
\usepackage{tabularx,booktabs,colortbl,array}
\usepackage{titlesec}
\usepackage[most]{tcolorbox}
\usepackage{enumitem}
\usepackage{fancyhdr}
\usepackage{hyperref}
\usepackage{xurl}
\usepackage{needspace}

% ---------- 色板（与 ochrelune-design 一致）----------
\definecolor{paper}{HTML}{FAF9F5}
\definecolor{oat}{HTML}{F0EEE6}
\definecolor{card}{HTML}{E8E6DC}
\definecolor{ink}{HTML}{141413}
\definecolor{ink2}{HTML}{5E5D59}
\definecolor{muted}{HTML}{87867F}
\definecolor{clay}{HTML}{D97757}
\definecolor{claydark}{HTML}{9F5038}
\definecolor{olive}{HTML}{788C5D}
\definecolor{rule}{HTML}{D9D5C8}
\colorlet{todoframe}{oat}
\def\todorule{0pt}
\ifdefined\ochreprint
  \definecolor{paper}{HTML}{FFFFFF}
  \definecolor{oat}{HTML}{FFFFFF}
  \definecolor{card}{HTML}{FFFFFF}
  \colorlet{todoframe}{rule}
  \def\todorule{0.6pt}
\fi
\pagecolor{paper}
\color{ink}

% ---------- 字体：标题寒蝉锦书宋 Pro Medium，大数字 Source Serif 4，正文霞鹜新晰黑 ----------
\setmainfont{LXGWNeoXiHei.ttf}[Path=fonts/,AutoFakeBold=1.6]
\setmonofont{Menlo}[Scale=0.88]
\setCJKmainfont{LXGWNeoXiHei.ttf}[Path=fonts/,AutoFakeBold=1.6]
\newfontfamily\serifLatin{ChillJinshuSongMedium.otf}[Path=fonts/]
\newCJKfontfamily\serifCJK{ChillJinshuSongMedium.otf}[Path=fonts/]
\newfontfamily\numfont{SourceSerif4.ttf}[Path=fonts/,Numbers=Lining]
\newcommand{\serif}{\serifLatin\serifCJK}
\xeCJKsetup{CJKecglue={\hskip 0.2em plus 0.1em}}
\linespread{1.42}
\setlength{\parindent}{0pt}
\setlength{\parskip}{0.55em}

\providecommand{\ochretitle}{Overview}
\providecommand{\ochrefooter}{\ochretitle}
\hypersetup{colorlinks=true,linkcolor=claydark,urlcolor=claydark,pdftitle={\ochretitle},pdfauthor={Fennelia}}

% ---------- 标题 ----------
\titleformat{\section}{\serif\fontsize{21}{28}\selectfont\color{ink}}{\numfont\color{clay}\thesection}{0.7em}{}[\vspace{2pt}{\color{rule}\titlerule[0.6pt]}]
\titlespacing*{\section}{0pt}{6pt}{12pt}
\titleformat{\subsection}{\serif\fontsize{14}{20}\selectfont\color{ink}}{}{0em}{}
\titlespacing*{\subsection}{0pt}{14pt}{4pt}

% ---------- 页脚 ----------
\pagestyle{fancy}
\fancyhf{}
\renewcommand{\headrulewidth}{0pt}
\fancyfoot[L]{\footnotesize\color{muted}\ochrefooter}
\fancyfoot[R]{\small\numfont\color{muted}\thepage}

% ---------- 表格 ----------
\newcolumntype{L}{>{\raggedright\arraybackslash}X}
\newcolumntype{P}[1]{>{\raggedright\arraybackslash}p{#1}}
\arrayrulecolor{rule}
\setlength{\aboverulesep}{0.5ex}
\setlength{\belowrulesep}{0.6ex}
\renewcommand{\arraystretch}{1.18}
\newcommand{\thd}[1]{{\footnotesize\color{muted}#1}}

% ---------- 组件 ----------
% todo：她要做或要决定的事（燕麦底 + 左侧陶土竖线）
\newtcolorbox{todo}[1][你要做的]{enhanced,breakable,colback=oat,colframe=todoframe,boxrule=\todorule,arc=3pt,
  left=12pt,right=12pt,top=8pt,bottom=8pt,
  borderline west={2.4pt}{0pt}{clay},
  title={\serif\color{claydark}#1},coltitle=claydark,colbacktitle=oat,attach title to upper,after title={\par\smallskip}}
% advice：我的建议（纸面底 + 细边框）
\newtcolorbox{advice}[1][我的建议]{enhanced,breakable,colback=paper,colframe=rule,boxrule=0.6pt,arc=3pt,
  left=12pt,right=12pt,top=8pt,bottom=8pt,
  title={\serif\color{ink}#1},attach title to upper,after title={\par\smallskip}}

% 元信息行：\meta{截止}{执行会话}{材料}
\newcommand{\metaline}[2]{\par\noindent\makebox[0.15\linewidth][l]{\small\color{muted}#1}\parbox[t]{0.85\linewidth}{\small\raggedright #2}\par\vspace{3pt}}
\newcommand{\meta}[3]{{\setlength{\parskip}{0pt}\metaline{截止}{#1}\metaline{执行会话}{#2}\metaline{材料}{#3}}\vspace{4pt}}

\setlist[itemize]{leftmargin=1.3em,itemsep=0.25em,topsep=0.3em,label={\color{clay}\textbullet}}
\setlist[enumerate]{leftmargin=1.5em,itemsep=0.3em,topsep=0.3em,label={\numfont\color{claydark}\arabic*.}}

% 封面大数字（一行最多四个）：\bignum{6}{天后答辩\\11-24 周二}
\newcommand{\bignum}[2]{\begin{minipage}[t]{0.23\linewidth}{\numfont\fontsize{40}{40}\selectfont\color{ink}#1}\par\vspace{4pt}{\small\color{ink2}#2}\end{minipage}}
\newcommand{\fpath}[1]{{\small\color{ink2}#1}}
\newcommand{\hairline}{{\color{rule}\rule{\linewidth}{0.6pt}}\par}

% 需要在主文件之前（或主文件第一行之前）定义好：
%   \def\ochretitle{文档标题}   \def\ochrefooter{主题 · 日期}（页脚不写署名）
% 打印版：再加一行 \def\ochreprint{}，纸面和卡片底改为纯白（用户自备米色纸，黑白打印机印底色会出淡墨）。
% 字体放在同目录 fonts/ 下（assets/fetch-fonts.sh 会下载）。
\documentclass[11pt]{article}
\usepackage[a4paper,top=22mm,bottom=24mm,left=22mm,right=22mm,headsep=8mm,footskip=12mm]{geometry}
\usepackage{fontspec}
\usepackage[CJKmath=true]{xeCJK}
\usepackage{xcolor}
\usepackage{tabularx,booktabs,colortbl,array}
\usepackage{titlesec}
\usepackage[most]{tcolorbox}
\usepackage{enumitem}
\usepackage{fancyhdr}
\usepackage{hyperref}
\usepackage{xurl}
\usepackage{needspace}

% ---------- 色板（与 ochrelune-design 一致）----------
\definecolor{paper}{HTML}{FAF9F5}
\definecolor{oat}{HTML}{F0EEE6}
\definecolor{card}{HTML}{E8E6DC}
\definecolor{ink}{HTML}{141413}
\definecolor{ink2}{HTML}{5E5D59}
\definecolor{muted}{HTML}{87867F}
\definecolor{clay}{HTML}{D97757}
\definecolor{claydark}{HTML}{9F5038}
\definecolor{olive}{HTML}{788C5D}
\definecolor{rule}{HTML}{D9D5C8}
\colorlet{todoframe}{oat}
\def\todorule{0pt}
\ifdefined\ochreprint
  \definecolor{paper}{HTML}{FFFFFF}
  \definecolor{oat}{HTML}{FFFFFF}
  \definecolor{card}{HTML}{FFFFFF}
  \colorlet{todoframe}{rule}
  \def\todorule{0.6pt}
\fi
\pagecolor{paper}
\color{ink}

% ---------- 字体：标题寒蝉锦书宋 Pro Medium，大数字 Source Serif 4，正文霞鹜新晰黑 ----------
\setmainfont{LXGWNeoXiHei.ttf}[Path=fonts/,AutoFakeBold=1.6]
\setmonofont{Menlo}[Scale=0.88]
\setCJKmainfont{LXGWNeoXiHei.ttf}[Path=fonts/,AutoFakeBold=1.6]
\newfontfamily\serifLatin{ChillJinshuSongMedium.otf}[Path=fonts/]
\newCJKfontfamily\serifCJK{ChillJinshuSongMedium.otf}[Path=fonts/]
\newfontfamily\numfont{SourceSerif4.ttf}[Path=fonts/,Numbers=Lining]
\newcommand{\serif}{\serifLatin\serifCJK}
\xeCJKsetup{CJKecglue={\hskip 0.2em plus 0.1em}}
\linespread{1.42}
\setlength{\parindent}{0pt}
\setlength{\parskip}{0.55em}

\providecommand{\ochretitle}{Overview}
\providecommand{\ochrefooter}{\ochretitle}
\hypersetup{colorlinks=true,linkcolor=claydark,urlcolor=claydark,pdftitle={\ochretitle},pdfauthor={Fennelia}}

% ---------- 标题 ----------
\titleformat{\section}{\serif\fontsize{21}{28}\selectfont\color{ink}}{\numfont\color{clay}\thesection}{0.7em}{}[\vspace{2pt}{\color{rule}\titlerule[0.6pt]}]
\titlespacing*{\section}{0pt}{6pt}{12pt}
\titleformat{\subsection}{\serif\fontsize{14}{20}\selectfont\color{ink}}{}{0em}{}
\titlespacing*{\subsection}{0pt}{14pt}{4pt}

% ---------- 页脚 ----------
\pagestyle{fancy}
\fancyhf{}
\renewcommand{\headrulewidth}{0pt}
\fancyfoot[L]{\footnotesize\color{muted}\ochrefooter}
\fancyfoot[R]{\small\numfont\color{muted}\thepage}

% ---------- 表格 ----------
\newcolumntype{L}{>{\raggedright\arraybackslash}X}
\newcolumntype{P}[1]{>{\raggedright\arraybackslash}p{#1}}
\arrayrulecolor{rule}
\setlength{\aboverulesep}{0.5ex}
\setlength{\belowrulesep}{0.6ex}
\renewcommand{\arraystretch}{1.18}
\newcommand{\thd}[1]{{\footnotesize\color{muted}#1}}

% ---------- 组件 ----------
% todo：她要做或要决定的事（燕麦底 + 左侧陶土竖线）
\newtcolorbox{todo}[1][你要做的]{enhanced,breakable,colback=oat,colframe=todoframe,boxrule=\todorule,arc=3pt,
  left=12pt,right=12pt,top=8pt,bottom=8pt,
  borderline west={2.4pt}{0pt}{clay},
  title={\serif\color{claydark}#1},coltitle=claydark,colbacktitle=oat,attach title to upper,after title={\par\smallskip}}
% advice：我的建议（纸面底 + 细边框）
\newtcolorbox{advice}[1][我的建议]{enhanced,breakable,colback=paper,colframe=rule,boxrule=0.6pt,arc=3pt,
  left=12pt,right=12pt,top=8pt,bottom=8pt,
  title={\serif\color{ink}#1},attach title to upper,after title={\par\smallskip}}

% 元信息行：\meta{截止}{执行会话}{材料}
\newcommand{\metaline}[2]{\par\noindent\makebox[0.15\linewidth][l]{\small\color{muted}#1}\parbox[t]{0.85\linewidth}{\small\raggedright #2}\par\vspace{3pt}}
\newcommand{\meta}[3]{{\setlength{\parskip}{0pt}\metaline{截止}{#1}\metaline{执行会话}{#2}\metaline{材料}{#3}}\vspace{4pt}}

\setlist[itemize]{leftmargin=1.3em,itemsep=0.25em,topsep=0.3em,label={\color{clay}\textbullet}}
\setlist[enumerate]{leftmargin=1.5em,itemsep=0.3em,topsep=0.3em,label={\numfont\color{claydark}\arabic*.}}

% 封面大数字（一行最多四个）：\bignum{6}{天后答辩\\11-24 周二}
\newcommand{\bignum}[2]{\begin{minipage}[t]{0.23\linewidth}{\numfont\fontsize{40}{40}\selectfont\color{ink}#1}\par\vspace{4pt}{\small\color{ink2}#2}\end{minipage}}
\newcommand{\fpath}[1]{{\small\color{ink2}#1}}
\newcommand{\hairline}{{\color{rule}\rule{\linewidth}{0.6pt}}\par}

% 需要在主文件之前（或主文件第一行之前）定义好：
%   \def\ochretitle{文档标题}   \def\ochrefooter{主题 · 日期}（页脚不写署名）
% 打印版：再加一行 \def\ochreprint{}，纸面和卡片底改为纯白（用户自备米色纸，黑白打印机印底色会出淡墨）。
% 字体放在同目录 fonts/ 下（assets/fetch-fonts.sh 会下载）。
\documentclass[11pt]{article}
\usepackage[a4paper,top=22mm,bottom=24mm,left=22mm,right=22mm,headsep=8mm,footskip=12mm]{geometry}
\usepackage{fontspec}
\usepackage[CJKmath=true]{xeCJK}
\usepackage{xcolor}
\usepackage{tabularx,booktabs,colortbl,array}
\usepackage{titlesec}
\usepackage[most]{tcolorbox}
\usepackage{enumitem}
\usepackage{fancyhdr}
\usepackage{hyperref}
\usepackage{xurl}
\usepackage{needspace}

% ---------- 色板（与 ochrelune-design 一致）----------
\definecolor{paper}{HTML}{FAF9F5}
\definecolor{oat}{HTML}{F0EEE6}
\definecolor{card}{HTML}{E8E6DC}
\definecolor{ink}{HTML}{141413}
\definecolor{ink2}{HTML}{5E5D59}
\definecolor{muted}{HTML}{87867F}
\definecolor{clay}{HTML}{D97757}
\definecolor{claydark}{HTML}{9F5038}
\definecolor{olive}{HTML}{788C5D}
\definecolor{rule}{HTML}{D9D5C8}
\colorlet{todoframe}{oat}
\def\todorule{0pt}
\ifdefined\ochreprint
  \definecolor{paper}{HTML}{FFFFFF}
  \definecolor{oat}{HTML}{FFFFFF}
  \definecolor{card}{HTML}{FFFFFF}
  \colorlet{todoframe}{rule}
  \def\todorule{0.6pt}
\fi
\pagecolor{paper}
\color{ink}

% ---------- 字体：标题寒蝉锦书宋 Pro Medium，大数字 Source Serif 4，正文霞鹜新晰黑 ----------
\setmainfont{LXGWNeoXiHei.ttf}[Path=fonts/,AutoFakeBold=1.6]
\setmonofont{Menlo}[Scale=0.88]
\setCJKmainfont{LXGWNeoXiHei.ttf}[Path=fonts/,AutoFakeBold=1.6]
\newfontfamily\serifLatin{ChillJinshuSongMedium.otf}[Path=fonts/]
\newCJKfontfamily\serifCJK{ChillJinshuSongMedium.otf}[Path=fonts/]
\newfontfamily\numfont{SourceSerif4.ttf}[Path=fonts/,Numbers=Lining]
\newcommand{\serif}{\serifLatin\serifCJK}
\xeCJKsetup{CJKecglue={\hskip 0.2em plus 0.1em}}
\linespread{1.42}
\setlength{\parindent}{0pt}
\setlength{\parskip}{0.55em}

\providecommand{\ochretitle}{Overview}
\providecommand{\ochrefooter}{\ochretitle}
\hypersetup{colorlinks=true,linkcolor=claydark,urlcolor=claydark,pdftitle={\ochretitle},pdfauthor={Fennelia}}

% ---------- 标题 ----------
\titleformat{\section}{\serif\fontsize{21}{28}\selectfont\color{ink}}{\numfont\color{clay}\thesection}{0.7em}{}[\vspace{2pt}{\color{rule}\titlerule[0.6pt]}]
\titlespacing*{\section}{0pt}{6pt}{12pt}
\titleformat{\subsection}{\serif\fontsize{14}{20}\selectfont\color{ink}}{}{0em}{}
\titlespacing*{\subsection}{0pt}{14pt}{4pt}

% ---------- 页脚 ----------
\pagestyle{fancy}
\fancyhf{}
\renewcommand{\headrulewidth}{0pt}
\fancyfoot[L]{\footnotesize\color{muted}\ochrefooter}
\fancyfoot[R]{\small\numfont\color{muted}\thepage}

% ---------- 表格 ----------
\newcolumntype{L}{>{\raggedright\arraybackslash}X}
\newcolumntype{P}[1]{>{\raggedright\arraybackslash}p{#1}}
\arrayrulecolor{rule}
\setlength{\aboverulesep}{0.5ex}
\setlength{\belowrulesep}{0.6ex}
\renewcommand{\arraystretch}{1.18}
\newcommand{\thd}[1]{{\footnotesize\color{muted}#1}}

% ---------- 组件 ----------
% todo：她要做或要决定的事（燕麦底 + 左侧陶土竖线）
\newtcolorbox{todo}[1][你要做的]{enhanced,breakable,colback=oat,colframe=todoframe,boxrule=\todorule,arc=3pt,
  left=12pt,right=12pt,top=8pt,bottom=8pt,
  borderline west={2.4pt}{0pt}{clay},
  title={\serif\color{claydark}#1},coltitle=claydark,colbacktitle=oat,attach title to upper,after title={\par\smallskip}}
% advice：我的建议（纸面底 + 细边框）
\newtcolorbox{advice}[1][我的建议]{enhanced,breakable,colback=paper,colframe=rule,boxrule=0.6pt,arc=3pt,
  left=12pt,right=12pt,top=8pt,bottom=8pt,
  title={\serif\color{ink}#1},attach title to upper,after title={\par\smallskip}}

% 元信息行：\meta{截止}{执行会话}{材料}
\newcommand{\metaline}[2]{\par\noindent\makebox[0.15\linewidth][l]{\small\color{muted}#1}\parbox[t]{0.85\linewidth}{\small\raggedright #2}\par\vspace{3pt}}
\newcommand{\meta}[3]{{\setlength{\parskip}{0pt}\metaline{截止}{#1}\metaline{执行会话}{#2}\metaline{材料}{#3}}\vspace{4pt}}

\setlist[itemize]{leftmargin=1.3em,itemsep=0.25em,topsep=0.3em,label={\color{clay}\textbullet}}
\setlist[enumerate]{leftmargin=1.5em,itemsep=0.3em,topsep=0.3em,label={\numfont\color{claydark}\arabic*.}}

% 封面大数字（一行最多四个）：\bignum{6}{天后答辩\\11-24 周二}
\newcommand{\bignum}[2]{\begin{minipage}[t]{0.23\linewidth}{\numfont\fontsize{40}{40}\selectfont\color{ink}#1}\par\vspace{4pt}{\small\color{ink2}#2}\end{minipage}}
\newcommand{\fpath}[1]{{\small\color{ink2}#1}}
\newcommand{\hairline}{{\color{rule}\rule{\linewidth}{0.6pt}}\par}

% 需要在主文件之前（或主文件第一行之前）定义好：
%   \def\ochretitle{文档标题}   \def\ochrefooter{主题 · 日期}（页脚不写署名）
% 打印版：再加一行 \def\ochreprint{}，纸面和卡片底改为纯白（用户自备米色纸，黑白打印机印底色会出淡墨）。
% 字体放在同目录 fonts/ 下（assets/fetch-fonts.sh 会下载）。
\documentclass[11pt]{article}
\usepackage[a4paper,top=22mm,bottom=24mm,left=22mm,right=22mm,headsep=8mm,footskip=12mm]{geometry}
\usepackage{fontspec}
\usepackage[CJKmath=true]{xeCJK}
\usepackage{xcolor}
\usepackage{tabularx,booktabs,colortbl,array}
\usepackage{titlesec}
\usepackage[most]{tcolorbox}
\usepackage{enumitem}
\usepackage{fancyhdr}
\usepackage{hyperref}
\usepackage{xurl}
\usepackage{needspace}

% ---------- 色板（与 ochrelune-design 一致）----------
\definecolor{paper}{HTML}{FAF9F5}
\definecolor{oat}{HTML}{F0EEE6}
\definecolor{card}{HTML}{E8E6DC}
\definecolor{ink}{HTML}{141413}
\definecolor{ink2}{HTML}{5E5D59}
\definecolor{muted}{HTML}{87867F}
\definecolor{clay}{HTML}{D97757}
\definecolor{claydark}{HTML}{9F5038}
\definecolor{olive}{HTML}{788C5D}
\definecolor{rule}{HTML}{D9D5C8}
\colorlet{todoframe}{oat}
\def\todorule{0pt}
\ifdefined\ochreprint
  \definecolor{paper}{HTML}{FFFFFF}
  \definecolor{oat}{HTML}{FFFFFF}
  \definecolor{card}{HTML}{FFFFFF}
  \colorlet{todoframe}{rule}
  \def\todorule{0.6pt}
\fi
\pagecolor{paper}
\color{ink}

% ---------- 字体：标题寒蝉锦书宋 Pro Medium，大数字 Source Serif 4，正文霞鹜新晰黑 ----------
\setmainfont{LXGWNeoXiHei.ttf}[Path=fonts/,AutoFakeBold=1.6]
\setmonofont{Menlo}[Scale=0.88]
\setCJKmainfont{LXGWNeoXiHei.ttf}[Path=fonts/,AutoFakeBold=1.6]
\newfontfamily\serifLatin{ChillJinshuSongMedium.otf}[Path=fonts/]
\newCJKfontfamily\serifCJK{ChillJinshuSongMedium.otf}[Path=fonts/]
\newfontfamily\numfont{SourceSerif4.ttf}[Path=fonts/,Numbers=Lining]
\newcommand{\serif}{\serifLatin\serifCJK}
\xeCJKsetup{CJKecglue={\hskip 0.2em plus 0.1em}}
\linespread{1.42}
\setlength{\parindent}{0pt}
\setlength{\parskip}{0.55em}

\providecommand{\ochretitle}{Overview}
\providecommand{\ochrefooter}{\ochretitle}
\hypersetup{colorlinks=true,linkcolor=claydark,urlcolor=claydark,pdftitle={\ochretitle},pdfauthor={Fennelia}}

% ---------- 标题 ----------
\titleformat{\section}{\serif\fontsize{21}{28}\selectfont\color{ink}}{\numfont\color{clay}\thesection}{0.7em}{}[\vspace{2pt}{\color{rule}\titlerule[0.6pt]}]
\titlespacing*{\section}{0pt}{6pt}{12pt}
\titleformat{\subsection}{\serif\fontsize{14}{20}\selectfont\color{ink}}{}{0em}{}
\titlespacing*{\subsection}{0pt}{14pt}{4pt}

% ---------- 页脚 ----------
\pagestyle{fancy}
\fancyhf{}
\renewcommand{\headrulewidth}{0pt}
\fancyfoot[L]{\footnotesize\color{muted}\ochrefooter}
\fancyfoot[R]{\small\numfont\color{muted}\thepage}

% ---------- 表格 ----------
\newcolumntype{L}{>{\raggedright\arraybackslash}X}
\newcolumntype{P}[1]{>{\raggedright\arraybackslash}p{#1}}
\arrayrulecolor{rule}
\setlength{\aboverulesep}{0.5ex}
\setlength{\belowrulesep}{0.6ex}
\renewcommand{\arraystretch}{1.18}
\newcommand{\thd}[1]{{\footnotesize\color{muted}#1}}

% ---------- 组件 ----------
% todo：她要做或要决定的事（燕麦底 + 左侧陶土竖线）
\newtcolorbox{todo}[1][你要做的]{enhanced,breakable,colback=oat,colframe=todoframe,boxrule=\todorule,arc=3pt,
  left=12pt,right=12pt,top=8pt,bottom=8pt,
  borderline west={2.4pt}{0pt}{clay},
  title={\serif\color{claydark}#1},coltitle=claydark,colbacktitle=oat,attach title to upper,after title={\par\smallskip}}
% advice：我的建议（纸面底 + 细边框）
\newtcolorbox{advice}[1][我的建议]{enhanced,breakable,colback=paper,colframe=rule,boxrule=0.6pt,arc=3pt,
  left=12pt,right=12pt,top=8pt,bottom=8pt,
  title={\serif\color{ink}#1},attach title to upper,after title={\par\smallskip}}

% 元信息行：\meta{截止}{执行会话}{材料}
\newcommand{\metaline}[2]{\par\noindent\makebox[0.15\linewidth][l]{\small\color{muted}#1}\parbox[t]{0.85\linewidth}{\small\raggedright #2}\par\vspace{3pt}}
\newcommand{\meta}[3]{{\setlength{\parskip}{0pt}\metaline{截止}{#1}\metaline{执行会话}{#2}\metaline{材料}{#3}}\vspace{4pt}}

\setlist[itemize]{leftmargin=1.3em,itemsep=0.25em,topsep=0.3em,label={\color{clay}\textbullet}}
\setlist[enumerate]{leftmargin=1.5em,itemsep=0.3em,topsep=0.3em,label={\numfont\color{claydark}\arabic*.}}

% 封面大数字（一行最多四个）：\bignum{6}{天后答辩\\11-24 周二}
\newcommand{\bignum}[2]{\begin{minipage}[t]{0.23\linewidth}{\numfont\fontsize{40}{40}\selectfont\color{ink}#1}\par\vspace{4pt}{\small\color{ink2}#2}\end{minipage}}
\newcommand{\fpath}[1]{{\small\color{ink2}#1}}
\newcommand{\hairline}{{\color{rule}\rule{\linewidth}{0.6pt}}\par}
